\section{Lexical Analysis}\label{sec:lexical}

\subsection{Scanner design}

The lexical layer is written from scratch -- no parser generator, no \texttt{lex}. It runs
in two passes over the raw statement.

\textbf{Pass one -- scanning.} \texttt{core/text.py} splits the input into typed character
runs (word, number, punctuation, separator, space, newline) using the five regular
expressions in Table~\ref{tab:regex}. Each run keeps its byte offsets so every later
artefact -- token, parse error, repair suggestion -- can point back at the exact source
span the user typed.

\textbf{Pass two -- lexing.} \texttt{core/lexer.py} consumes those runs and resolves each
word against the lexicon, producing tokens carrying a terminal class, a canonical form, a
donor-language label and a gloss.

\begin{table}[htbp]
\centering
\caption{Regular expressions driving the scanner, exported from \texttt{core/text.py}.}
\label{tab:regex}
\begin{tabular}{@{}l p{0.40\linewidth} p{0.32\linewidth}@{}}
\toprule
\textbf{Span kind} & \textbf{Regular expression} & \textbf{Matches} \\
\midrule
\texttt{WORD} & \ttfamily\scriptsize (\allowbreak{}?\allowbreak{}:\allowbreak{}[\allowbreak{}\textasciicircum{}\allowbreak{}\textbackslash{}\allowbreak{}W\allowbreak{}\textbackslash{}\allowbreak{}d\allowbreak{}\_\allowbreak{}]\allowbreak{}|\allowbreak{}[\allowbreak{}\textbackslash{}\allowbreak{}u\allowbreak{}0\allowbreak{}3\allowbreak{}0\allowbreak{}0\allowbreak{}-\allowbreak{}\textbackslash{}\allowbreak{}u\allowbreak{}0\allowbreak{}3\allowbreak{}6\allowbreak{}f\allowbreak{}]\allowbreak{})\allowbreak{}+\allowbreak{}(\allowbreak{}?\allowbreak{}:\allowbreak{}[\allowbreak{}'\allowbreak{}\textbackslash{}\allowbreak{}u\allowbreak{}2\allowbreak{}0\allowbreak{}1\allowbreak{}9\allowbreak{}-\allowbreak{}]\allowbreak{}(\allowbreak{}?\allowbreak{}:\allowbreak{}[\allowbreak{}\textasciicircum{}\allowbreak{}\textbackslash{}\allowbreak{}W\allowbreak{}\textbackslash{}\allowbreak{}d\allowbreak{}\_\allowbreak{}]\allowbreak{}|\allowbreak{}[\allowbreak{}\textbackslash{}\allowbreak{}u\allowbreak{}0\allowbreak{}3\allowbreak{}0\allowbreak{}0\allowbreak{}-\allowbreak{}\textbackslash{}\allowbreak{}u\allowbreak{}0\allowbreak{}3\allowbreak{}6\allowbreak{}f\allowbreak{}]\allowbreak{})\allowbreak{}+\allowbreak{})\allowbreak{}* & Letters and combining accents, with internal apostrophes or hyphens \\
\addlinespace
\texttt{NUMBER} & \ttfamily\scriptsize \textbackslash{}\allowbreak{}d\allowbreak{}+ & A run of digits \\
\addlinespace
\texttt{SPACE} & \ttfamily\scriptsize [\allowbreak{}\textasciicircum{}\allowbreak{}\textbackslash{}\allowbreak{}S\allowbreak{}\textbackslash{}\allowbreak{}r\allowbreak{}\textbackslash{}\allowbreak{}n\allowbreak{}]\allowbreak{}+ & Horizontal whitespace only \\
\addlinespace
\texttt{NEWLINE} & \ttfamily\scriptsize [\allowbreak{}\textbackslash{}\allowbreak{}r\allowbreak{}\textbackslash{}\allowbreak{}n\allowbreak{}]\allowbreak{}+ & One or more line breaks \\
\addlinespace
\texttt{PUNCTUATION} & \ttfamily\scriptsize [\allowbreak{}.\allowbreak{},\allowbreak{}!\allowbreak{}?\allowbreak{};\allowbreak{}:\allowbreak{}\textbackslash{}\allowbreak{}"\allowbreak{}(\allowbreak{})\allowbreak{}«\allowbreak{}»\allowbreak{}\textbackslash{}\allowbreak{}u\allowbreak{}2\allowbreak{}0\allowbreak{}1\allowbreak{}8\allowbreak{}\textbackslash{}\allowbreak{}u\allowbreak{}2\allowbreak{}0\allowbreak{}1\allowbreak{}9\allowbreak{}\textbackslash{}\allowbreak{}u\allowbreak{}2\allowbreak{}0\allowbreak{}1\allowbreak{}c\allowbreak{}\textbackslash{}\allowbreak{}u\allowbreak{}2\allowbreak{}0\allowbreak{}1\allowbreak{}d\allowbreak{}…\allowbreak{}]\allowbreak{}+ & Sentence and clause marks, including guillemets \\
\addlinespace
\bottomrule
\end{tabular}
\end{table}


\subsection{Word regex in detail}

The word pattern is the one that does real linguistic work:

\begin{lstlisting}[language=Python, numbers=none, basicstyle=\ttfamily\footnotesize]
_WORD_RE = (?:[^\W\d_]|[\u0300-\u036f])+
           (?:['\u2019-](?:[^\W\d_]|[\u0300-\u036f])+)*
\end{lstlisting}

\begin{itemize}[leftmargin=1.4em, itemsep=0.2em]
  \item \term{[\^{}\textbackslash W\textbackslash d\_]} matches any Unicode letter while
        excluding digits and underscore, so \fw{ça}, \fw{où} and \fw{tchop} all scan as
        single words without an explicit accent list.
  \item \term{[\textbackslash u0300-\textbackslash u036f]} admits \emph{combining} marks,
        so a decomposed \fw{e + U+0301} scans identically to a precomposed \fw{é}. This is
        essential for typed input, where the two encodings are visually identical.
  \item The trailing group keeps elisions and hyphenated compounds whole:
        \fw{j'ai}, \fw{n'est}, \fw{est-ce}, \fw{quelqu'un}. Both the ASCII apostrophe and
        the typographic \fw{U+2019} are accepted, because phones produce the latter.
\end{itemize}

\subsection{Token categories}

\begin{tabular}{@{}lrrl@{}}
\toprule
\textbf{Terminal} & \textbf{Entries} & \textbf{Slang} & \textbf{Example} \\
\midrule
VERB & 290 & 97 & \texttt{go} \\
NOUN & 210 & 192 & \texttt{combi} \\
ADV & 35 & 13 & \texttt{même} \\
ADJ & 33 & 23 & \texttt{cher} \\
DET & 30 & 1 & \texttt{les} \\
PRON & 23 & 1 & \texttt{on} \\
PREP & 22 & 0 & \texttt{au} \\
INTJ & 19 & 19 & \texttt{ashia} \\
CONJ & 14 & 0 & \texttt{qu'il} \\
FORMULA & 7 & 7 & \texttt{on dit quoi} \\
NEG & 5 & 0 & \texttt{ne} \\
\midrule
\textbf{Total} & \textbf{688} & \textbf{353} & \\
\bottomrule
\end{tabular}


The \TerminalCount\ terminals cover the categories the rubric asks for and three that the
data forced on us:

\begin{itemize}[leftmargin=1.4em, itemsep=0.2em]
  \item \term{NOUN}, \term{VERB}, \term{ADJ}, \term{ADV}, \term{DET}, \term{PRON},
        \term{PREP}, \term{CONJ} -- the conventional word classes.
  \item \term{NEG} is split out from \term{ADV} because Francanglais negation is
        positionally constrained in a way ordinary adverbs are not, and folding the two
        together introduced grammar conflicts.
  \item \term{INTJ} covers the discourse particles (\fw{ehh}, \fw{ashia}, \fw{wah}) that
        open and close a very large share of real utterances.
  \item \term{FORMULA} covers fixed multiword expressions that are lexically frozen --
        they are matched as one token even though they span several words.
  \item \term{SEP} is emitted for \term{,} \term{;} \term{:}. These are genuine clause
        boundaries in speech, so the parser needs to see them; sentence-final
        \term{.\;!\;?} remain trivia.
\end{itemize}

\subsection{Slang and code-mixed tokens}

Every entry records the donor language it came from, which is what lets the analyzer
report code-mixing quantitatively rather than impressionistically.

\begin{tabular}{@{}lr@{}}
\toprule
\textbf{Donor language} & \textbf{Token entries} \\
\midrule
french & 425 \\
uncertain & 132 \\
english & 50 \\
pidgin & 23 \\
mixed & 14 \\
duala & 14 \\
beti & 11 \\
bassa & 6 \\
bamileke & 5 \\
fulfulde & 3 \\
hausa & 2 \\
bulu & 1 \\
german & 1 \\
lingala & 1 \\
spanish & 1 \\
\bottomrule
\end{tabular}


A statement is reported as code-mixed when its tokens resolve to more than one donor
label; the reported evidence is the list of tokens and labels themselves. Slang is not a
separate terminal -- a slang noun is still a \term{NOUN} -- because slang is a fact about
\emph{provenance}, not about syntax. Treating it otherwise would have duplicated the whole
grammar.

\subsection{Accent and case normalisation}

Two normalisations are applied, and the distinction between them matters:

\begin{description}[leftmargin=1.6em, style=nextline, font=\ttfamily]
  \item[canonicalize\_word] NFC normalise, then casefold. This is the identity used for
        storage and display: \fw{Taxi}, \fw{TAXI} and \fw{taxi} are one entry.
  \item[fold\_word] NFD decompose, strip combining marks, recompose, casefold. This is the
        \emph{fallback} identity: \fw{bénéf} and \fw{benef} reach the same entry even
        though the user typed neither form exactly.
\end{description}

Folding is deliberately a fallback rather than the primary key, because in French it
destroys real contrasts. Lookup therefore proceeds in strict priority order:

\begin{enumerate}[leftmargin=1.4em, itemsep=0.1em]
  \item exact multiword match, \item folded multiword match,
  \item exact single-word match, \item folded single-word match,
  \item otherwise \term{UNKNOWN}.
\end{enumerate}

Because exact matches win, the four attested homograph pairs survive intact:

\begin{center}
\begin{tabular}{@{}llll@{}}
\toprule
\textbf{Form A} & \textbf{Class} & \textbf{Form B} & \textbf{Class} \\
\midrule
\fw{a}    & \term{VERB} (\textit{has})      & \fw{à}    & \term{PREP} (\textit{to, at}) \\
\fw{la}   & \term{DET} (\textit{the})       & \fw{là}   & \term{ADV} (\textit{there}) \\
\fw{ou}   & \term{CONJ} (\textit{or})       & \fw{où}   & \term{ADV} (\textit{where}) \\
\fw{mais} & \term{CONJ} (\textit{but})      & \fw{maïs} & \term{NOUN} (\textit{maize}) \\
\bottomrule
\end{tabular}
\end{center}

When a folded lookup resolves to a form that has a homograph, the analyzer emits a
homograph hint naming both readings instead of choosing one silently -- the user is told
that \fw{la} was read as a determiner and that \fw{là} exists.

\subsection{Multiword expressions}

\LexiconMultiword\ entries span more than one word (\fw{est-ce que}, \fw{on va faire
comment}). The lexer indexes them by first word, attempts the longest match first, and
refuses to match across a punctuation or separator boundary -- so a comma between two words
prevents them from being glued into one token. This rule is covered by a dedicated
regression test.

\subsection{Unknown words}

An unrecognised word is not an error and is not discarded. It is emitted as an
\term{UNKNOWN} token with its source span intact, which produces three useful effects: the
parser fails at a precise location rather than vaguely; the suggester can offer near
matches computed over \emph{folded} forms, so an accent slip still finds its target; and
the reviewer gets a concrete list of vocabulary the lexicon is missing.

The \LexiconCorpus\ corpus-attested entries, each with the statement identifiers that
justify it, are listed in Appendix~\ref{app:lexicon}.

\section{Token Frequency and Variation}\label{sec:frequency}

Three frequency views are produced for every analysis, and they answer different questions.

\textbf{Raw surface forms} show what people actually typed or said, accents, casing and
spelling variants included.

\begin{tabular}{@{}lr@{}}
\toprule
\textbf{Raw form} & \textbf{Count} \\
\midrule
\texttt{,} & 20 \\
\texttt{le} & 15 \\
\texttt{au} & 10 \\
\texttt{Le} & 9 \\
\texttt{est} & 6 \\
\texttt{je} & 6 \\
\texttt{la} & 6 \\
\texttt{pas} & 6 \\
\texttt{kwatt} & 5 \\
\texttt{on} & 4 \\
\texttt{va} & 4 \\
\texttt{Combi} & 3 \\
\texttt{a} & 3 \\
\texttt{go} & 3 \\
\texttt{les} & 3 \\
\texttt{ne} & 3 \\
\texttt{piol} & 3 \\
\texttt{à} & 3 \\
\texttt{Ashia} & 2 \\
\texttt{Je} & 2 \\
\bottomrule
\end{tabular}


\textbf{Canonical forms} collapse casing and Unicode normalisation differences. Comparing
this table against the raw one measures orthographic variation directly: a canonical form
reached by several distinct surface spellings is a documented variant cluster.

\begin{tabular}{@{}lr@{}}
\toprule
\textbf{Canonical form} & \textbf{Count} \\
\midrule
\texttt{le} & 24 \\
\texttt{,} & 20 \\
\texttt{au} & 10 \\
\texttt{je} & 8 \\
\texttt{est} & 6 \\
\texttt{la} & 6 \\
\texttt{on} & 6 \\
\texttt{pas} & 6 \\
\texttt{kwatt} & 5 \\
\texttt{combi} & 4 \\
\texttt{va} & 4 \\
\texttt{a} & 3 \\
\texttt{go} & 3 \\
\texttt{les} & 3 \\
\texttt{ne} & 3 \\
\texttt{piol} & 3 \\
\texttt{à} & 3 \\
\texttt{ashia} & 2 \\
\texttt{avec} & 2 \\
\texttt{bring} & 2 \\
\bottomrule
\end{tabular}


\textbf{Terminal classes} show the grammatical shape of the data, which is what predicts
whether a grammar will parse it.

\begin{tabular}{@{}lr@{}}
\toprule
\textbf{Terminal} & \textbf{Count} \\
\midrule
\texttt{NOUN} & 60 \\
\texttt{VERB} & 58 \\
\texttt{DET} & 40 \\
\texttt{PREP} & 25 \\
\texttt{SEP} & 20 \\
\texttt{PRON} & 19 \\
\texttt{ADV} & 14 \\
\texttt{NEG} & 9 \\
\texttt{ADJ} & 6 \\
\texttt{CONJ} & 4 \\
\texttt{INTJ} & 2 \\
\texttt{FORMULA} & 1 \\
\bottomrule
\end{tabular}


\begin{keypoint}[What the variation tells us]
Variation in this data is not noise to be normalised away. The gap between the raw and
canonical tables is a measurement of how unstable written Francanglais orthography is,
and it is precisely the reason the lexer needs a folded fallback tier rather than a single
exact-match dictionary.
\end{keypoint}
